\section{Controlled Experimental Setup}
\label{app:controlled-setup}

The controlled suite gives DECAF an identifiable target before we move to natural images. Every experiment uses exact factor interventions in 3D Shapes, shared protocol randomness across branches, and two architectures. The experiments differ in the behavior they are designed to activate.

\subsection{Dataset and factors}

3D Shapes contains $480{,}000$ rendered scenes generated from six factors: floor color, wall color, object color, object size, object shape, and object orientation~\citep{burgess2018understanding}. The Cartesian product is fully enumerated. This structure gives us exact counterfactuals: to intervene on one factor, we change its index and hold the other five indices fixed.

We use binary prediction tasks derived from the factors. The base benchmark contains five tasks:
\begin{itemize}[leftmargin=1.4em,itemsep=2pt,topsep=3pt]
    \item \textbf{Object color:} a binary split of the object-color values.
    \item \textbf{Wall color:} a binary split of the wall-color values.
    \item \textbf{Object shape:} the two shapes that are common to all model-response supports.
    \item \textbf{Color--shape XOR:} the exclusive-or of binary object color and shape.
    \item \textbf{Context gate:} a color decision whose active factor depends on wall context.
\end{itemize}
For each task we train ResNet-18 and a small ViT with three random seeds. The resulting 30 models are crossed with all six factors, producing 180 model--factor units.

\subsection{Clean endpoint pairs}

For a factual scene $\mathbf{x}^{+}$ and target factor $k$, the counterfactual $\mathbf{x}^{-}$ changes only factor $k$. Binary factors are flipped. Multivalued color, size, and orientation factors use fixed involutive maps so that applying the map twice returns the original value. We use two independent maps for multivalued factors and average their measurements.

The clean endpoint effect is computed on 8,192 held-out factual--counterfactual pairs. The main threshold classifies a pair as active when the absolute target-score difference exceeds the experiment-specific $\varepsilon$. The endpoint classification is a property of the model and pair, not of the task label.

\subsection{Primary reveal protocol}
\label{app:cmmr}

The main controlled protocol is covariance-matched Gaussian reveal. Let $\boldsymbol\mu$ and $\widehat{\boldsymbol\Sigma}$ be the empirical mean and covariance of the image distribution. For $\alpha\in[0,1]$,
\begin{equation}
\mathbf{x}_{\alpha}
=
\boldsymbol\mu
+
\sqrt{\alpha}(\mathbf{x}-\boldsymbol\mu)
+
\sqrt{1-\alpha}\,\boldsymbol\eta,
\qquad
\boldsymbol\eta\sim\mathcal N(\mathbf{0},\widehat{\boldsymbol\Sigma}).
\label{eq:cmmr-path}
\end{equation}
The factual and counterfactual branch share the same $\boldsymbol\eta$. At $\alpha=1$ the path reaches the clean input; at $\alpha=0$ the state is independent of the individual sample under the fitted Gaussian reference. The covariance match makes the stimulus second-order neutral in whitened data coordinates. DECAF does not depend on this particular path; it only consumes the paired scores that the path produces.

The base benchmark uses 4,096 dynamic factual pairs, three noise seeds, two counterfactual maps, and 21 points along a trace-matched covariance family. The primary CMMR endpoint is accompanied by a trace-matched pixel-Gaussian endpoint. A diagonal covariance and a trace-normalized covariance-power path provide held-out geometries.

\subsection{Alternative protocols}

The covariance interpolation is
\[
\boldsymbol\Gamma_{\lambda}
=(1-\lambda)\widehat{\boldsymbol\Sigma}
+
\lambda\tau\mathbf I,
\qquad
\tau=\frac{\operatorname{tr}(\widehat{\boldsymbol\Sigma})}{D},
\]
where $D$ is the pixel dimension. The two endpoints use data covariance and trace-matched isotropic pixel covariance. The diagonal held-out path uses $\operatorname{diag}(\widehat{\boldsymbol\Sigma})$. The power path raises covariance eigenvalues to a power and renormalizes the trace. These paths test whether a semantic conclusion depends on one second-order geometry.

We use the protocol family as an audit, not as a learned worst-case score. The main paper reports CMMR-based DECAF and uses the alternatives to characterize transfer.

\subsection{Evidence validation through training trajectories}
\label{app:e-setup}

The evidence experiment trains shape classifiers in an environment where wall color is strongly correlated with the label. We use two training correlations, $0.95$ and $0.99$, two architectures, two seeds, and save checkpoints throughout training. This creates eight trajectories and 52 selected checkpoints.

The external target is shortcut-reversal vulnerability,
\[
V_{\mathrm{rev}}
=
\operatorname{Acc}(p_{\mathrm{test}}=0.95)
-
\operatorname{Acc}(p_{\mathrm{test}}=0.05).
\]
A shortcut-dominant checkpoint performs well when wall color remains aligned and fails when the correlation reverses. A shape-dominant checkpoint is stable. We compare $V_{\mathrm{rev}}$ with the evidence margin $E_{\mathrm{wall}}-E_{\mathrm{shape}}$.

\subsection{Fragility validation}
\label{app:f-setup}

The fragility experiment uses a binary object-shape task and treats floor color as the candidate endpoint-null factor. We train robust, neutral, and fragile variants for each architecture and three seeds, producing 18 models.

The neutral variant uses clean training only. The robust variant is encouraged to keep its output stable under floor-color counterfactuals at intermediate states. The fragile variant is encouraged to respond to floor color away from the endpoint while preserving the clean shape task. The intended comparison is restricted to models for which clean accuracy remains high and floor color remains endpoint-null.

The behavioral target is the prediction-change rate under held-out floor-color interventions at intermediate states. We also measure randomized-floor stability and transfer across covariance geometries.

\subsection{Contradiction validation}
\label{app:c-setup}

The contradiction benchmark uses two binary variables: object color $A$ and wall context $G$. The factual endpoint is evaluated in the context $G=1$, where all three tasks share the clean rule $Y=A$. Under the swapped context $G=0$:
\[
\text{Direct: }Y=A,
\qquad
\text{Gate: }Y=H,
\qquad
\text{Invert: }Y=1-A,
\]
where $H$ is object shape. Direct preserves the object-color effect. Gate removes it. Invert reverses it.

We train 30 models: three tasks, two architectures, and five seeds. Balanced validation accuracy is at least $0.9987$. A binary symmetric context channel swaps the wall context with probability $\eta\in[0,0.5]$. All intermediate images remain in the support of the generated dataset. Two independent wall-color involutions test map transfer.

The label-level outcomes are preserve, collapse, and swap. The key contrast is between Gate and Invert: both reduce aligned evidence, but only Invert should create endpoint-opposed mass.

\subsection{Statistical summaries}

The base benchmark uses stratified bootstrap over task, architecture, and seed. The evidence experiment bootstraps complete trajectories. The fragility and contradiction experiments report seed-level variation and architecture-stratified summaries. Bootstrap intervals use 500 repetitions unless stated otherwise. We treat the model or training trajectory as the main statistical unit rather than counting every image as independent evidence.
